% Part: proof-theory
% Chapter: natural-deduction
% Section: translation-G2i

\documentclass[../../../include/open-logic-section]{subfiles}

\begin{document}

\olfileid[id]{pt}{nat}{tgi}

\olsection{Menerjemahkan dari \Log{G2i} ke \Log{N2i}}

Anteseden sekuen dalam \Log{G2i} adalah multihimpunan~$\Gamma$ dari
!!{formula}, sedangkan anteseden sekuen dalam \Log{N2i} adalah
himpunan~$\Gamma'$ dari formula berlabel dengan sifat bahwa setiap label
hanya muncul satu kali. Kita katakan bahwa suatu himpunan~$\Gamma'$ dari
!!{formula} berlabel \emph{bersesuaian dengan} multihimpunan~$\Gamma$
jika, untuk setiap !!{formula}~$!A$, banyaknya !!{element} dalam
$\Gamma'$ yang berbentuk $x:!A$ kurang dari atau sama dengan
multiplisitas~$!A$ (yakni banyaknya salinan~$!A$) dalam~$\Gamma$.

\begin{prop}\ollabel{prop:G2i-to-N2i}
Jika $\Log{G2i}+\Cut \Proves \Gamma \Sequent \Delta$, maka
$\Log{N2i} \Proves \Gamma' \Sequent !C$, dengan
$!C\ident\lfalse$ jika $\Delta$ kosong dan
$\Delta=\{!A\}$ jika tidak, serta $\Gamma'$ bersesuaian
dengan~$\Gamma$.
\end{prop}

\begin{proof}
  Kita menunjukkan bahwa jika $\pi$ adalah !!a{proof} atas
  $\Gamma \Sequent \Delta$ dalam $\Log{G2i}+\Cut$, maka terdapat
  suatu~$\Gamma'$ yang bersesuaian dengan~$\Gamma$ dan
  !!a{proof}~$\delta$ atas~$\Gamma'\Sequent !C$ dalam~\Log{N2i}.
  Kita melakukannya melalui induksi pada $n=\pheight{\pi}$.

  Jika $n=0$, maka $\pi$ merupakan aksioma. Jika $\pi$ adalah aksioma
  $\lfalse\Sequent\quad$, maka $\delta$ adalah aksioma
  $x:\lfalse\Sequent\lfalse$; yakni,
  $\Gamma'=\{x:\lfalse\}$ dan $!C\ident\lfalse$. Jika $\pi$ merupakan
  aksioma berbentuk $!A\Sequent !A$, maka $\delta$ adalah
  $x:!A\Sequent !A$.

  Jika $n>0$, $\pi$ berakhir pada suatu aturan~$R$. Hipotesis induksi
  menjamin adanya !!{proof} atas sekuen-sekuen yang bersesuaian dengan
  premis aturan tersebut dalam~\Log{N2i}. Kita dapat mengasumsikan bahwa
  semua label pelepasan dalam bukti-bukti itu berbeda, dan bahwa suatu
  label $x$ hanya muncul di atas inferensi yang melepaskannya, jika
  label itu dilepaskan sama sekali (dengan mengganti label bila perlu).
  Lalu kita menunjukkan cara menggabungkan !!{proof} tersebut
  untuk menghasilkan bukti atas $\Gamma'\Sequent !C$ yang sesuai.
  Kita membedakan kasus menurut~$R$.

  \begin{enumerate}
    \item Jika $R$ adalah \RightR{\Weakening}, $\pi$ berakhir pada
    \[
    \AxiomC{}
    \RightLabel{$\pi_1$}
    \Deduce$\Gamma \fCenter$
    \RightLabel{\RightR{\Weakening}}
    \UnaryInf$\Gamma \fCenter !A$
    \DisplayProof
    \]
    Hipotesis induksi yang diterapkan pada $\pi_1$ menghasilkan
    !!a{proof} atas~$\Gamma'\Sequent\lfalse$. Kita menerapkan
    $\FalseInt$ untuk memperoleh !!a{proof} atas
    $\Gamma'\Sequent !A$.

    \item Jika $R$ adalah \LeftR{\Weakening}, $\pi$ berakhir pada
    \[
    \AxiomC{}
    \RightLabel{$\pi_1$}
    \Deduce$\Gamma \fCenter \Delta$
    \RightLabel{\LeftR{\Weakening}}
    \UnaryInf$!A, \Gamma \fCenter \Delta$
    \DisplayProof
    \]
    Hipotesis induksi yang diterapkan pada $\pi_1$ menghasilkan
    !!a{proof} atas~$\Gamma'\Sequent !C$, yang kita ambil
    sebagai~$\delta$. Perhatikan bahwa jika $\Gamma'$ bersesuaian
    dengan~$\Gamma$, maka $\Gamma'$ juga bersesuaian dengan
    $!A,\Gamma$: banyaknya $x:!A$ dalam $\Gamma'$ tidak melebihi
    banyaknya salinan~$!A$ dalam~$\Gamma$, sehingga juga tidak melebihi
    banyaknya salinan $!A$ dalam $\Gamma\cup\{!A\}$.

    \item Jika $R$ adalah \LeftR{\Contraction}, $\pi$ berakhir pada
    \[
    \AxiomC{}
    \RightLabel{$\pi_1$}
    \Deduce$!A, !A, \Gamma \fCenter \Delta$
    \RightLabel{\LeftR{\Contraction}}
    \UnaryInf$!A, \Gamma \fCenter \Delta$
    \DisplayProof
    \]
    Hipotesis induksi yang diterapkan pada $\pi_1$ menghasilkan
    !!a{proof} atas~$\Pi,\Gamma'\Sequent !C$, dengan
    $\Pi\subseteq\{x:!A,y:!A\}$. Jika
    $\Pi=\{x:!A,y:!A\}$, ganti semua !!{formula} berlabel $y:!A$
    dalam~$\delta$ dengan $x:!A$. Karena $x$ muncul dalam konteks
    sekuen akhir~$\delta$, kita dapat mengasumsikan bahwa tidak ada
    inferensi dalam~$\delta$ yang melepaskan !!a{formula} berbentuk
    $x:!B$; yakni, tidak ada $x:!B$ di sebelah kiri sekuen
    dalam~$\delta$ yang aktif. Akibatnya, hasilnya tetap merupakan
    !!a{proof} dalam~\Log{N2i}.

    \item Jika $R$ adalah $\RightR{\lif}$, $\pi$ berakhir pada
    \[
    \AxiomC{}
    \RightLabel{$\pi_1$}
    \Deduce$!A, \Gamma \fCenter !B$
    \RightLabel{\RightR{\lif}}
    \UnaryInf$\Gamma \fCenter !A \lif !B$
    \DisplayProof
    \]
    Berdasarkan hipotesis induksi, terdapat !!a{proof}~$\delta_1$
    atas $x:!A,\Gamma'\Sequent !B$ atau
    $\Gamma'\Sequent !B$, dengan $\Gamma'$ bersesuaian
    dengan~$\Gamma$. Kita mengambil $\delta$ sebagai $\delta_1$ yang
    diikuti satu penerapan \Intro{\lif}.

    \item Jika $R$ adalah \LeftR{\lif}, $\pi$ berakhir pada
    \[
    \AxiomC{}
    \RightLabel{$\pi_1$}
    \Deduce$\Gamma_1 \fCenter !A$
    \AxiomC{}
    \RightLabel{$\pi_2$}
    \Deduce$!B, \Gamma_2 \fCenter \Delta$
    \RightLabel{\LeftR{\lif}}
    \BinaryInf$!A \lif !B, \Gamma_1, \Gamma_2 \fCenter \Delta$
    \DisplayProof
    \]
    Hipotesis induksi menghasilkan !!{proof}~$\delta_1$ atas
    $\Gamma_1'\Sequent !A$ serta $\delta_2$ atas
    $x:!B,\Gamma_2'\Sequent !C$ atau atas
    $\Gamma_2'\Sequent !C$. Jika $\delta_2$ berbentuk yang terakhir,
    $\delta_2$ dapat berperan sebagai~$\delta$. Jika tidak, kita dapat
    mengganti semua formula dan label pelepasan dalam~$\delta_1$
    untuk memastikan bahwa tidak ada label yang muncul sekaligus dalam
    $\delta_1$ dan~$\delta_2$. Maka
    $\Gamma_1'\cup\Gamma_2'$ bersesuaian dengan
    $\Gamma_1\cup\Gamma_2$. Karena $x:!B$ muncul dalam sekuen akhir
    $\delta_2$, $x:!B$ tidak dapat aktif dalam inferensi apa pun
    berlabel pelepasan~$x$ dalam~$\delta_2$. Ganti setiap aksioma
    $x:!B\Sequent !B$ dalam~$\delta_2$ dengan !!{proof}~$\delta_3$
    \[
    \Axiom$x: !A \lif !B \fCenter !A \lif !B$
    \AxiomC{}
    \RightLabel{$\delta_1$}
    \Deduce$\Gamma_1' \fCenter !A$
    \RightLabel{\Elim{\lif}}
    \BinaryInf$x: !A \lif !B, \Gamma_1' \fCenter !B$
    \DisplayProof
    \]
    dan ganti setiap kemunculan $x:!B$ dalam~$\delta_2$ dengan
    $x:!A\lif !B$. Hasil
    $\Subst{\delta_2}{\delta_3}{x:!B}$ merupakan !!a{proof} atas
    $x:!A\lif !B,\Gamma_1',\Gamma_2'\Sequent !C$.

    \item Jika $R$ adalah \RightR{\land}, $\pi$ berakhir pada
    \[
    \AxiomC{}
    \RightLabel{$\pi_1$}
    \Deduce$\Gamma_1 \fCenter !A$
    \AxiomC{}
    \RightLabel{$\pi_2$}
    \Deduce$\Gamma_2 \fCenter !B$
    \RightLabel{\RightR{\land}}
    \BinaryInf$\Gamma_1, \Gamma_2 \fCenter !A \land !B$
    \DisplayProof
    \]
    Hipotesis induksi menghasilkan bukti~$\delta_1$ atas
    $\Gamma_1'\Sequent !A$ dan $\delta_2$ atas
    $\Gamma_2'\Sequent !B$. Dengan mengganti semua formula dan label
    pelepasan dalam~$\delta_2$, kita dapat memastikan bahwa tidak ada
    label yang muncul sekaligus dalam $\delta_1$ dan $\delta_2$.
    Maka $\Gamma_1'\cup\Gamma_2'$ bersesuaian dengan
    $\Gamma_1\cup\Gamma_2$. Kita memperoleh !!a{proof}~$\delta$
    dengan menerapkan $\Intro{\land}$ pada sekuen akhir $\delta_1$
    dan~$\delta_2$.

    \item Jika $R$ adalah $\LeftR{\land}$, $\pi$ berakhir, misalnya,
    pada
    \[
    \AxiomC{}
    \RightLabel{$\pi_1$}
    \Deduce$!A, \Gamma \fCenter \Delta$
    \RightLabel{\LeftR{\land}}
    \UnaryInf$!A \land !B, \Gamma \fCenter \Delta$
    \DisplayProof
    \]
    Berdasarkan hipotesis induksi, terdapat !!a{proof}~$\delta_1$
    atas $x:!A,\Gamma'\Sequent !C$ atau
    $\Gamma'\Sequent !C$, dengan $\Gamma'$ bersesuaian
    dengan~$\Gamma$. Dalam kasus kedua, kita dapat mengambil
    $\delta_1$ sebagai~$\delta$. Dalam kasus pertama, perhatikan bahwa
    karena $x:!A$ muncul dalam sekuen akhir, $x:!A$ tidak aktif dalam
    inferensi mana pun berlabel pelepasan~$x$ dalam~$\delta_1$.
    Ganti setiap aksioma $x:!A\Sequent !A$ dengan !!{proof}~$\delta_2$
    \[
    \Axiom$x: !A \land !B \fCenter !A \land !B$
    \RightLabel{\Elim{\land}}
    \UnaryInf$x: !A \land !B \fCenter !A$
    \DisplayProof
    \]
    dan setiap kemunculan $x:!A$ dengan $x:!A\land !B$; yakni,
    tinjau $\Subst{\delta_1}{\delta_2}{x:!A}$. Ini merupakan
    !!a{proof} atas
    $x:!A\land !B,\Gamma'\Sequent !C$.

    \item Jika $R$ adalah \Cut, $\pi$ berakhir pada
    \[
    \AxiomC{}
    \RightLabel{$\pi_1$}
    \Deduce$\Gamma_1 \fCenter !A$
    \AxiomC{}
    \RightLabel{$\pi_2$}
    \Deduce$!A, \Gamma_2 \fCenter \Delta$
    \RightLabel{\Cut}
    \BinaryInf$\Gamma_1, \Gamma_2 \fCenter \Delta$
    \DisplayProof
    \]
    Hipotesis induksi menghasilkan bukti~$\delta_1$ atas
    $\Gamma_1'\Sequent !A$ dan $\delta_2$ atas
    $x:!A,\Gamma_2'\Sequent !C$ atau atas
    $\Gamma_2'\Sequent !C$. Dalam kasus kedua, ambil $\delta$ sebagai
    $\delta_2$. Jika tidak, ganti setiap aksioma
    $x:!A\Sequent !A$ dalam $\delta_2$ dengan $\delta_1$ untuk
    memperoleh~$\delta$ sebagai !!{proof}
    $\Subst{\delta_2}{\delta_1}{x:!A}$ atas
    $\Gamma_1',\Gamma_2'\Sequent !C$.
  \end{enumerate}
\end{proof}

\end{document}
